\section{Introduction}
\label{sec:introduction}

% STRUCTURE ONLY. Written after production; no conclusions before results.
%
% The argument this section must make, in order:
%
% 1. Survival prediction models are routinely reported with a concordance
%    index, and a concordance index is a property of a *ranking*. It is
%    invariant to any strictly monotone transformation of the risk score, so a
%    model whose predicted probabilities are uniformly wrong attains the same
%    value as one that is exactly right. State this as the formal fact it is:
%    it is what makes the question non-vacuous.
%
% 2. That discrimination and calibration can disagree is KNOWN. Say so here,
%    plainly and early, with citations. Do not stage it as a discovery. See
%    \citet{austin2020graphical}, \citet{sonabend2022chacking},
%    \citet{birolo2025beyond}, \citet{lillelund2026stop}.
%
% 3. The gap: all of that work evaluates against censored observations, which
%    do not contain the true conditional survival function. So the size of the
%    error a conventional metric permits has not been measured -- only its
%    existence demonstrated.
%
% 4. The contribution: evaluation against an analytically known S(t | x),
%    which permits the integrated squared distance between the predicted and
%    the true individual survival curve directly.
%
% 5. The question, stated as a measurement: when conventional metrics still
%    describe a model as adequate, how far has the estimated individual
%    survival distribution moved from the truth?
%
% 6. Contributions list. Be precise about which are novel and which reproduce.
